\section{Related work}
\label{sec:related-work}

\subsection{Runs and the speed at which deposits leave}

The theory of deposit runs starts from the observation that a demand-deposit contract gives depositors liquidity and also admits an equilibrium in which all of them withdraw at once \citep{diamond1983}. \citet{goldsteinpauzner2005} make fundamentals decide which equilibrium occurs, which lets them compute the ex ante probability of a panic-based run and relate it to the contract. Both works explain why a run can happen; neither says how fast deposits leave, which is the quantity a liquidity stress test needs. Depositor-level evidence supplies it in part: \citet{iyerpuri2012} follow the depositors of a bank that faced a run and find that uninsured depositors are the most likely to run, that deposit insurance helps only partially, that relationships with the bank mitigate the run and that social networks matter. \citet{cipriani2024} trace the deposit flows of March 2023 from high-frequency interbank payments, identify twenty-two banks that suffered a run, and find that large institutional depositors drove it, that coordination mattered, and that banks that survived a run did so by borrowing new funds and raising deposit rates rather than by selling liquid securities. % voice-ok: withdraw is the depositors action in a bank run, not a version narrative

The funding side of the balance sheet has its own literature. \citet{brunnermeierpedersen2009} model the link between the market liquidity of an asset and the funding liquidity of traders, and show that under certain conditions margins are destabilising and the two reinforce each other in liquidity spirals; \citet{drehmannnikolaou2013} define funding liquidity as the ability to settle obligations with immediacy and propose a measure of its risk built on bids at central bank auctions; \citet{gortonmetrick2012} argue that the panic of 2007 to 2008 was a run on the repo market, with securitised banking at its centre, and document it with the credit spreads of securitised bonds and repo rates; and \citet{acharyamora2015} document that at the onset of the 2007 to 2009 crisis aggregate deposit inflows weakened and loan-to-deposit shortfalls widened, most at banks with greater undrawn commitments. These works explain why funding dries up and where it is fragile; none gives the speed of outflow at which a stated buffer is exhausted. \citet{tirole2011} reviews illiquidity and explains why liquidity cannot easily be apprehended through a single statistic, which is the reason the reverse stress below is stated for a survival period under a stated scenario and not as a general liquidity measure. For 2023, \citet{jiang2024} mark the assets of United States banks to market after the rise in interest rates and examine the role of unrealised losses and uninsured depositors in the runs; they report an average decline of about 10\% in the value of bank assets, about 2 trillion US dollars in aggregate. \citet{cookson2026} study social media as a catalyst of the run on the failed bank; their working-paper abstract reports that banks with high pre-run exposure to Twitter lost more stock market value during the run, and we have not checked the figures of the published version. In the abstract of the working paper under its first title, ``Banking on Uninsured Deposits'',\citet{drechsler2023} show that a run equilibrium among uninsured depositors is absent at low interest rates and appears when rates rise, because the deposit franchise then dominates the value of the bank, so that the liquidity risk of a bank increases with interest rates. What these studies imply for a stress test is that the speed of a run depends on the type of depositor, on the information environment and on the value of the securities that back the deposits, and that it cannot be calibrated once and for all. A reverse stress that reports the speed at which a target is reached is a statement about distance, not about likelihood.

\subsection{Regulatory and supervisory liquidity stress}

The Basel Committee's principles for liquidity risk management ask a bank to stress-test short-term and protracted, institution-specific and market-wide scenarios, individually and in combination \citep{bcbs144}. The liquidity coverage ratio fixes one combined thirty-day scenario with prescribed run-off rates \citep{bcbs238}, and the net stable funding ratio addresses structural funding over a year \citep{bcbs295}. The Committee's principles for stress testing of 2009, since superseded by its later principles, set out the governance, design and implementation of stress-testing programmes and the expectations of supervisors \citep{bcbs155}. The supervisory stress test of the European Central Bank in 2019 specified smooth scenarios: open-maturity deposits run off at a constant daily compounding rate on relevant days, contractual flows roll off at stated rates when they mature, and the survival period is the first day on which the net liquidity position turns negative \citep{ecblist2019,ecbannex2019}. The Committee's report on the 2023 turmoil gives a table of observed deposit outflows over a number of days, from 8\% in 14 days at IndyMac to 85\% in two days at the bank that failed, and records that the assumptions of the liquidity coverage ratio, 3\% to 40\% over thirty days, lie far from the fastest of them \citep{bcbs555}. The Federal Reserve's review of the failed bank found that it repeatedly failed its own internal liquidity stress test and that the deposit assumptions of that test relied on peer benchmarks that were not comparable \citep{fedsvb2023}.

These documents fix scenarios and report the survival they give. They do not give the multiplier at which a stated scenario reaches a stated survival period, and they do not say what an exercise should report when no admissible multiplier does.

\subsection{Liquidity stress-testing models}

Models that implement liquidity stress tests differ in how they combine the run scenario with the rest of the balance sheet. \citet{vandenend2009} proposes a liquidity stress-tester that includes the first- and second-round effects of shocks, the reactions of heterogeneous banks and reputation effects, simulated by Monte Carlo and applied to Dutch banks; \citet{puhr2012} present a framework for system-wide, balance-sheet-based liquidity stress tests with a module that simulates bank-run scenarios, a module for maturity transformation and rollover risk, and a link between liquidity and solvency risks; and \citet{cont2020} propose a structural framework for the joint stress testing of solvency and liquidity and define Liquidity at Risk as the liquidity resources that a stress scenario requires. \citet{borio2014} appraise macro stress testing and conclude that, with current methods, it is ill-suited as an early-warning device and more useful as a tool of crisis management; \citet{schuermann2014} sets out a framework for bank stress testing: why it became the supervisory tool in the crisis, how it is designed and executed, and how its results should be disclosed. The engine used here, the \texttt{riskvalidation} package, implements the ladder of the supervisory exercise (contractual and behavioural ladders, counterbalancing capacity with haircuts and a monetisation schedule, net liquidity position, survival period) and the reverse stress module that this paper studies. We do not compare it with the models above because they answer a different question: each takes a scenario as an input, and the paper's object is the inversion of that input for one outcome.

\subsection{Reverse stress tests}

Reverse stress tests were developed mostly for market, credit and macroeconomic risk. \citet{glasserman2015} estimate the most likely market-factor scenarios given that portfolio losses exceed a threshold, using an empirical-likelihood estimator, with justification for multivariate normal, Laplace and $t$ factor models; we read the abstract and not the full text. \citet{flood2015} select multidimensional shock scenarios by a grid search of points of arbitrary severity drawn from a plausible joint distribution, which they regard as a middle ground between traditional and reverse stress testing; \citet{breuercsiszar2013} look for the worst expected portfolio value among the scenarios that are sufficiently plausible, plausibility being measured by an entropic distance from a reference distribution; and \citet{mcneilsmith2012} build convex scenario sets for the risk factors of an insurer by half-space trimming and take the least solvent likely event. \citet{grundke2018} propose a quantitative method for maturity-transforming banks that finds the macroeconomic scenarios that lead to default and attaches probabilities to them, with copulas and extreme value theory. All of these search a space of risk factors and need a density or a plausibility measure to choose among scenarios. We characterise the works of this section from their abstracts or records and not from the full texts. Book-length treatments exist: \citet{bellini2017} devotes a chapter to reverse stress testing in a text on stress testing and risk integration, \citet{gann2021} discusses reverse stress testing in banks, and \citet{almehed2021} wrote a chapter titled reverse liquidity stress testing. We could not read the text of the last three, so whether any of them defines a multiplier on run-off rates, solves it as a threshold or states a refusal rule is not established here and stays unverified; the claim of this paper is limited to a defined solver, its measured behaviour and its theory.

The proposed method differs from the market-factor approaches in what is searched. They seek a most plausible point on a loss surface, which needs a density on the factors. The liquidity multiplier is one-dimensional and monotone by construction, so no plausibility measure is needed, and the difficulty moves to the integer-valued outcome, the points at which the survival period is undefined, and the refusal of targets that no admissible scenario reaches.
